\section{Experiments}

We organize the experiments around three questions: Can \method{} deliver a useful accuracy-time trade-off for full-materialization TNC? Do the gains actually come from internal contraction rather than output write-back? And how do approximation strength and individual mechanisms affect the result?

\subsection{Experimental setup}

We report four representative classical networks: Ring-8~\cite{bondy2008graph}, Tree-8~\cite{diestel2017graph}, Grid $3\times 3$~\cite{levin2007trg,orus2014practical}, and Random-8~\cite{erdos1959random,gilbert1959random}. Here ``8'' denotes the number of nodes and $3\times 3$ denotes the two-dimensional grid size. Each node has one open leg, the physical dimension is fixed to $3$, and the bond dimension is fixed to $4$. We use \textsf{\method-L1/L2/L3} to denote different approximation strengths; unless otherwise stated, \textsf{\method} refers to the medium-strength default preset. We compare exact contraction (\textsf{exact-TNC}), a fixed-rank baseline, and \method{}. The fixed-rank baseline uses rank $4$ for all interfaces, serving as a clear small-rank fast-approximation reference. Additional details about network generation, tolerances, and output blocking are given in \Cref{app:exp-details}.

\subsection{Main results}

\Cref{tab:main-core} reports the core comparison across the four networks. The fixed-rank baseline is consistently very fast, but also severely distorted: its relative error ranges from $0.65$ to $0.98$ on all four cases. In contrast, the default \method{} configuration keeps the relative error below $4.81\times 10^{-3}$ in every case while still achieving speedups from $13.98\times$ to $84.05\times$ over \textsf{exact-TNC}. The most discriminative case is Grid $3\times 3$: exact contraction takes $37.20$ seconds, whereas \method{} reduces the total time to $0.4426$ seconds with only $4.08\times 10^{-3}$ relative error. Across the four networks, \method{} attains a geometric-mean speedup of $42.53\times$ with average relative error $2.68\times 10^{-3}$.

\begin{table*}[t]
\centering
\footnotesize
\caption{Main results on four representative networks. \method{} provides a strong accuracy-time trade-off compared with exact contraction and a fixed-rank approximation baseline.}
\label{tab:main-core}
\setlength{\tabcolsep}{6pt}
\begin{tabular}{lccccc}
\toprule
Case & exact-TNC time (s) & fixed-rank rel. error & fixed-rank speedup & \method\ rel. error & \method\ speedup \\
\midrule
Ring-8 & 3.0072 & 0.8058 & 123.39$\times$ & 1.83$\times 10^{-3}$ & 56.31$\times$ \\
Tree-8 & 1.1068 & 0.6524 & 57.69$\times$ & 4.81$\times 10^{-3}$ & 49.45$\times$ \\
Grid $3\times 3$ & 37.1995 & 0.9789 & 566.07$\times$ & 4.08$\times 10^{-3}$ & 84.05$\times$ \\
Random-8 & 0.7326 & 0.9283 & 21.72$\times$ & $<10^{-12}$ & 13.98$\times$ \\
\midrule
Summary & \na & \na & \na & avg.\ $=2.68\times 10^{-3}$ & gmean $=42.53\times$ \\
\bottomrule
\end{tabular}
\end{table*}

\subsection{Where do the gains come from?}

We next examine whether the end-to-end gains are truly contraction-side. \Cref{tab:stage-breakdown-core} provides clear evidence: across the four main-result networks, the default \method{} configuration has average ratio $t_{\mathrm{contract}}/t_{\mathrm{total}} = 0.9992$, while its average contraction speedup and total speedup are $52.07\times$ and $52.03\times$, respectively. The two numbers are nearly identical. Hence the end-to-end acceleration is explained almost entirely by internal contraction rather than by output emission artifacts. For Grid $3\times 3$, the contraction speedup is $84.36\times$ and the total speedup is $84.33\times$, showing essentially no practical gap.

\begin{table}[t]
\centering
\small
\caption{Stage-time decomposition. The end-to-end gains of \method{} are almost entirely explained by contraction-side acceleration rather than output write-back.}
\label{tab:stage-breakdown-core}
\begin{adjustbox}{max width=\columnwidth}
\begin{tabular}{lcccc}
\toprule
Case & $t_{\mathrm{contract}}/t_{\mathrm{total}}$ & $t_{\mathrm{emit}}/t_{\mathrm{total}}$ & contraction speedup & total speedup \\
\midrule
Ring-8 & 0.9993 & 0.000694 & 56.37$\times$ & 56.33$\times$ \\
Tree-8 & 0.9985 & 0.001454 & 50.65$\times$ & 50.58$\times$ \\
Grid $3\times 3$ & 0.9997 & 0.000333 & 84.36$\times$ & 84.33$\times$ \\
Random-8 & 0.9993 & 0.000722 & 16.92$\times$ & 16.91$\times$ \\
\midrule
Average & 0.9992 & 0.000801 & 52.07$\times$ & 52.03$\times$ \\
\bottomrule
\end{tabular}
\end{adjustbox}
\end{table}

At the mechanism level, we use Grid $3\times 3$ as the most discriminative case and isolate two key components, as shown in \Cref{tab:grid-mechanism}. Removing implicit sketch-based merging leaves the error essentially unchanged but makes the total time $6.97\times$ worse than the default. This indicates that implicit sketching lowers internal merge cost without sacrificing accuracy. Removing output blocking increases the total time by $2.40\times$ and also raises the error from $4.08\times 10^{-3}$ to $1.06\times 10^{-2}$. For this harder network, implicit sketch-based merging is the dominant and most stable source of efficiency, while output blocking also provides a meaningful benefit.

\begin{table}[t]
\centering
\small
\caption{Mechanism evidence on Grid $3\times 3$. Implicit sketch-based merging is the dominant efficiency source, while output blocking also matters.}
\label{tab:grid-mechanism}
\begin{adjustbox}{max width=\columnwidth}
\begin{tabular}{lccc}
\toprule
Setting & rel. error & total time (s) & time ratio vs.\ default \\
\midrule
\method\ (default) & 4.08$\times 10^{-3}$ & 0.4426 & 1.00 \\
no-blockwise & 1.06$\times 10^{-2}$ & 1.0635 & 2.40 \\
no-implicit & 4.08$\times 10^{-3}$ & 3.0838 & 6.97 \\
\bottomrule
\end{tabular}
\end{adjustbox}
\end{table}

\subsection{Approximation-strength control}

We finally test whether approximation strength $\tau$ provides stable control over error. \Cref{tab:strength-sensitivity-core} shows a clear monotonic trend from \textsf{L1} to \textsf{L2} to \textsf{L3}. On Ring-8, the relative error increases from $<10^{-12}$ to $1.83\times 10^{-3}$ and then to $1.55\times 10^{-2}$. On Grid $3\times 3$, it increases from $<10^{-12}$ to $4.08\times 10^{-3}$ and then to $1.81\times 10^{-2}$. On Random-8, \textsf{L1} and \textsf{L2} are effectively exact, while \textsf{L3} rises to $1.26\times 10^{-2}$. The approximation level thus acts as a practical error knob for \method{}.

\begin{table}[t]
\centering
\small
\caption{Sensitivity to approximation strength. Stronger approximation leads to larger relative error in a consistent way.}
\label{tab:strength-sensitivity-core}
\begin{adjustbox}{max width=\columnwidth}
\begin{tabular}{lccc}
\toprule
Case & L1 rel. error & L2 rel. error & L3 rel. error \\
\midrule
Ring-8 & $<10^{-12}$ & 1.83$\times 10^{-3}$ & 1.55$\times 10^{-2}$ \\
Grid $3\times 3$ & $<10^{-12}$ & 4.08$\times 10^{-3}$ & 1.81$\times 10^{-2}$ \\
Random-8 & $<10^{-12}$ & $<10^{-12}$ & 1.26$\times 10^{-2}$ \\
\midrule
Average rel. error & $<10^{-12}$ & 1.97$\times 10^{-3}$ & 1.54$\times 10^{-2}$ \\
\bottomrule
\end{tabular}
\end{adjustbox}
\end{table}
